% TOP_PROBLEM: {"id":30007661,"problem_number":"LOCAL-30007661","title":"Reliable poverty estimates without fresh surveys","category_id":21,"category":{"id":21,"name":"economics","display_name":"Economics","description":"Open economic theory statements, published research agendas, and proposed scoped specifications; question provenance and historical priority are qualified per record.","slug":"economics","order_index":21,"created_at":"2026-10-08T00:00:00Z"},"status":"research_question","external_url":null,"legacy_ids":[],"published":false,"statement":"In the explicitly specified setting: How can poverty be estimated reliably between household surveys when economic shocks, changing survey design, and structural change invalidate past prediction relationships? The mathematical statement above fixes the target, assumptions and scope of this catalogue formulation.","economics":{"original_id":150,"field":"Development and poverty","kind":"identification","formalization_basis":"adapted_research_question","source_status":"explicit_open","priority_status":"earliest_found","priority_note":"Earliest directly inspected qualifying passage in this audit; earlier origins were not established. A review or research-agenda statement does not imply original priority.","earliest_verified_reference":{"authors":["Paul Corral","Henry Stemmler","Andrés Ham","Jaime Fernandez","Leonardo Lucchetti","Peter Lanjouw"],"title":"Bridging the survey gap: a call for innovative methods to predict poverty","year":2025,"url":"https://blogs.worldbank.org/en/opendata/bridging-the-survey-gap--a-call-for-innovative-methods-to-predic","doi":null,"locator":"Sections \"The challenge of predicting poverty\" and \"A competition to strengthen poverty prediction\", final three questions","evidence_summary":"The research authors explicitly ask which methods survive limited data and economic structural shifts and identify poverty reliably."},"current_reference":{"authors":["Paul Corral","Henry Stemmler","Andrés Ham","Jaime Fernandez","Leonardo Lucchetti","Peter Lanjouw"],"title":"Bridging the survey gap: a call for innovative methods to predict poverty","year":2025,"url":"https://blogs.worldbank.org/en/opendata/bridging-the-survey-gap--a-call-for-innovative-methods-to-predic","doi":null,"locator":"Sections \"The challenge of predicting poverty\" and \"A competition to strengthen poverty prediction\", final three questions","evidence_summary":"The research authors explicitly ask which methods survive limited data and economic structural shifts and identify poverty reliably."},"formalization_note":"The cited passage establishes a research direction, not this exact mathematical statement. The notation, population, policy menu, assumptions, and estimands are an original adaptation for this catalogue; no universal unsolved theorem or source priority is claimed."},"statusQualification":"Published question or research agenda verified at the cited locator. The mathematical specification is adapted for this catalogue and includes additional assumptions; source publication does not establish first-ever priority or certify this exact model remains unsolved. The cited passage establishes a research direction, not this exact mathematical statement. The notation, population, policy menu, assumptions, and estimands are an original adaptation for this catalogue; no universal unsolved theorem or source priority is claimed.","statusReviewedAt":"2026-10-08"}
% FIRST_POSTED: 2026-10-08
% ENTRY_KIND: statement_only
% !TeX program = lualatex
\documentclass[11pt]{article}
\usepackage{fontspec}
\usepackage[a4paper,margin=25mm]{geometry}
\usepackage{amsmath,amssymb,mathtools}
\usepackage{xurl}
\usepackage[hidelinks]{hyperref}
\newcommand{\sref}[2]{\hyperlink{source:#1:#2}{[S#2]}}
\setlength{\emergencystretch}{3em}
\begin{document}
% problemId:30007661
\section*{Reliable poverty estimates without fresh surveys}\label{30007661}
\subsection{Definitions and mathematical statement}
Fix two household-survey dates, \(0\) and \(1\), and a poverty threshold \(z>0\) expressed in comparable real-consumption units. At date \(0\), observe consumption \(C_0\) and predictors \(X_0\); at date \(1\), observe predictors \(X_1\) but not consumption \(C_1\). The target is \(H_1=\Pr(C_1<z)\), together with household poverty probabilities. Let \(\mathcal P\) be a declared class of joint distributions allowing specified changes in survey design, prices, predictor composition, and the conditional law of consumption. Construct an estimator \(\widehat H_1\) and interval \(I_1\) with valid stated coverage over \(\mathcal P\), or show which assumptions are necessary for such validity. Unrestricted conditional distribution shift makes the target unidentified; no method can solve that case solely by predictive accuracy at date \(0\). Compare candidate approaches on held-out surveys and deliberately simulated shifts, using both aggregate poverty error and household classification loss. Quantify which small new consumption samples or administrative measurements most improve robustness. The research authors explicitly pose reliability under structural shifts as an agenda. The distribution class, coverage criterion, and data-acquisition problem here are our adapted formal specification.

\par\textbf{Formulation and scope.} The cited passage establishes a research direction, not this exact mathematical statement. The notation, population, policy menu, assumptions, and estimands are an original adaptation for this catalogue; no universal unsolved theorem or source priority is claimed.

\par\textbf{Open-question provenance.} An explicit question or research agenda was verified at the source locator. This does not by itself certify that every later version remains unresolved \sref{30007661}{1}.

\par\textbf{Historical priority.} Earliest directly inspected qualifying passage in this audit; earlier origins were not established. A review or research-agenda statement does not imply original priority.

\subsection{Short English statement}
In the explicitly specified setting: How can poverty be estimated reliably between household surveys when economic shocks, changing survey design, and structural change invalidate past prediction relationships? The mathematical statement above fixes the target, assumptions and scope of this catalogue formulation.

\subsection{Sources}
\begin{itemize}\raggedright
\item[\textnormal{[S1]}]\hypertarget{source:30007661:1}{} Paul Corral, Henry Stemmler, Andrés Ham, Jaime Fernandez, Leonardo Lucchetti, Peter Lanjouw. 2025. Bridging the survey gap: a call for innovative methods to predict poverty. Sections "The challenge of predicting poverty" and "A competition to strengthen poverty prediction", final three questions.
\url{https://blogs.worldbank.org/en/opendata/bridging-the-survey-gap--a-call-for-innovative-methods-to-predic}.
{\small\textit{Earliest verified question/agenda reference; historical priority qualified below. The research authors explicitly ask which methods survive limited data and economic structural shifts and identify poverty reliably.}}
\end{itemize}
\end{document}
